\documentclass{tufte-handout}

%\geometry{showframe}% for debugging purposes -- displays the margins

\usepackage{amsmath}
\usepackage{listings}
\usepackage{color}

% Set up the images/graphics package
\usepackage{graphicx}
\setkeys{Gin}{width=\linewidth,totalheight=\textheight,keepaspectratio}
\graphicspath{{graphics/}}

\title{Tokenomics for LokuBox \\
\large The pluggable provenance tracking filesystem}
\author[Gary Mawdsley]{Gary Mawdsley CTO/CEO Lockular Limited}
\date{May 8th 2024\thanks{Work in progress}}  % if the \date{} command is left out, the current date will be used

% The following package makes prettier tables.  We're all about the communication!
\usepackage{booktabs}

% The units package provides nice, non-stacked fractions and better spacing
% for units.
\usepackage{units}

% The fancyvrb package lets us customize the formatting of verbatim
% environments.  We use a slightly smaller font.
\usepackage{fancyvrb}
\fvset{fontsize=\normalsize}

% Small sections of multiple columns
\usepackage{multicol}

% Provides paragraphs of dummy text
\usepackage{lipsum}
\usepackage{tabularx}
\usepackage{booktabs}

% These commands are used to pretty-print LaTeX commands
\newcommand{\doccmd}[1]{\texttt{\textbackslash#1}}% command name -- adds backslash automatically
\newcommand{\docopt}[1]{\ensuremath{\langle}\textrm{\textit{#1}}\ensuremath{\rangle}}% optional command argument
\newcommand{\docarg}[1]{\textrm{\textit{#1}}}% (required) command argument
\newenvironment{docspec}{\begin{quote}\noindent}{\end{quote}}% command specification environment
\newcommand{\docenv}[1]{\textsf{#1}}% environment name
\newcommand{\docpkg}[1]{\texttt{#1}}% package name
\newcommand{\doccls}[1]{\texttt{#1}}% document class name
\newcommand{\docclsopt}[1]{\texttt{#1}}% document class option name

\usepackage{titlesec} % Make sure to include this package

% Custom appendix command
\newcommand{\customappendix}{
    \clearpage % Ensure a page break
    \appendix % Mark the beginning of appendices
    \pagestyle{empty} % Optional: Removes header/footer if desired
    \titleformat{\section}[display]{\normalfont\Large\bfseries}{\appendixname~\thesection}{0pt}{\Large}
    \titlespacing*{\section}{0pt}{50pt}{40pt}
    \section*{Appendices} % This adds the title "Appendices"
    \addcontentsline{toc}{section}{Appendices} % Optional: Adds "Appendices" to the table of contents
}

% Define colors
\definecolor{codegreen}{rgb}{0,0.6,0}
\definecolor{codegray}{rgb}{0.5,0.5,0.5}
\definecolor{codepurple}{rgb}{0.58,0,0.82}
\definecolor{backcolour}{rgb}{0.95,0.95,0.92}

% Setup the listings package
\lstset{
    backgroundcolor=\color{backcolour},   
    commentstyle=\color{codegreen},
    keywordstyle=\color{magenta},
    numberstyle=\tiny\color{codegray},
    stringstyle=\color{codepurple},
    basicstyle=\ttfamily\footnotesize,
    breakatwhitespace=false,         
    breaklines=true,                 
    captionpos=b,                    
    keepspaces=true,                 
    numbers=left,                    
    numbersep=5pt,                  
    showspaces=false,                
    showstringspaces=false,
    showtabs=false,                  
    tabsize=2
}

\begin{document}

\section*{Creating a Tokenomics Model for LokuBox Service}

Creating a tokenomics model for a LokuBox service involves several key components. Here's a high-level overview of what the tokenomics model must consider:

\subsection*{1. Token Supply}
\begin{itemize}
    \item \textbf{Total Supply}: Define the maximum number of tokens that will ever exist.
    \item \textbf{Initial Supply}: The number of tokens available at launch.
\end{itemize}

\subsection*{2. Token Distribution}
\begin{itemize}
    \item \textbf{Founders and Team}: Percentage allocated to the founding team.
    \item \textbf{Advisors}: Tokens set aside for advisors.
    \item \textbf{Investors}: Tokens allocated for early investors or through an Initial Coin Offering (ICO).
    \item \textbf{Community and Ecosystem}: Tokens reserved for community incentives, partnerships, and ecosystem development.
    \item \textbf{Reserves}: Tokens held in reserve for future use.
\end{itemize}

\subsection*{3. Token Utility}
\begin{itemize}
    \item \textbf{Access to Services}: Tokens required to access or use the blockchain service.
    \item \textbf{Governance}: Tokens used for voting on protocol changes or governance decisions.
    \item \textbf{Staking}: Tokens that can be staked to earn rewards or participate in network validation.
    \item \textbf{Incentives}: Tokens used to incentivise certain behaviours, such as contributing to the network or providing liquidity.
\end{itemize}

\subsection*{4. Token Issuance and Inflation}
\begin{itemize}
    \item \textbf{Initial Distribution}: How tokens are initially distributed.
    \item \textbf{Inflation Rate}: If new tokens are minted over time, then the rate of inflation should be defined.
    \item \textbf{Burn Mechanism}: If applicable, describe how tokens are burned to reduce supply.
\end{itemize}

\subsection*{5. Economic Model}
\begin{itemize}
    \item \textbf{Revenue Model}: How the service generates revenue (e.g., transaction fees, subscription fees).
    \item \textbf{Token Flow}: How tokens move within the ecosystem (e.g., from users to service providers).
    \item \textbf{Market Dynamics}: Factors influencing token demand and supply.
\end{itemize}

\subsection*{6. Governance}
\begin{itemize}
    \item \textbf{Decision-Making Process}: How decisions are made within the ecosystem.
    \item \textbf{Voting Mechanism}: How token holders can vote on proposals.
\end{itemize}

\subsection*{7. Security and Compliance}
\begin{itemize}
    \item \textbf{Regulatory Compliance}: Ensuring the tokenomics model complies with relevant regulations.
    \item \textbf{Security Measures}: Steps taken to secure the token and the network.
\end{itemize}

\section*{Strawman Tokenomics Model for LokuBox}

Here's a tokenomics model for a LokuBox service:

\subsection*{Token Supply}
\begin{itemize}
    \item \textbf{Total Supply}: 1,000,000,000 tokens
    \item \textbf{Initial Supply}: 200,000,000 tokens
\end{itemize}

\subsection*{Token Distribution}
\begin{itemize}
    \item \textbf{Founders and Team}: 20\% (200,000,000 tokens)
    \item \textbf{Advisors}: 5\% (50,000,000 tokens)
    \item \textbf{Investors}: 15\% (150,000,000 tokens)
    \item \textbf{Community and Ecosystem}: 30\% (300,000,000 tokens)
    \item \textbf{Reserves}: 30\% (300,000,000 tokens)
\end{itemize}

\subsection*{Token Utility}
\begin{itemize}
    \item \textbf{Access to Services}: Tokens required to access premium features.
    \item \textbf{Governance}: Token holders can vote on protocol changes.
    \item \textbf{Staking}: Tokens can be staked to earn rewards.
    \item \textbf{Incentives}: Tokens used to reward users for contributing to the network.
\end{itemize}

\subsection*{Token Issuance and Inflation}
\begin{itemize}
    \item \textbf{Initial Distribution}: 200,000,000 tokens distributed at launch.
    \item \textbf{Inflation Rate}: 2\% annual inflation.
    \item \textbf{Burn Mechanism}: 1\% of transaction fees are burned.
\end{itemize}

\subsection*{Economic Model}
\begin{itemize}
    \item \textbf{Revenue Model}: Transaction fees, subscription fees for premium services.
    \item \textbf{Token Flow}: Users pay tokens to access services, service providers earn tokens.
    \item \textbf{Market Dynamics}: Demand driven by service usage, supply managed through staking and burning.
\end{itemize}

\subsection*{Governance}
\begin{itemize}
    \item \textbf{Decision-Making Process}: Proposals submitted by token holders, voted on by the community.
    \item \textbf{Voting Mechanism}: One token equals one vote.
\end{itemize}

\subsection*{Security and Compliance}
\begin{itemize}
    \item \textbf{Regulatory Compliance}: Adherence to local and international regulations.
    \item \textbf{Security Measures}: Regular audits, secure smart contract development.
\end{itemize}

Again this is a strawman framework and we'll tailor to fit the specific needs and goals of the LokuBox service as we review.

\section*{Inflation in a Tokenomics Model}

Inflation in a tokenomics model can serve several purposes:

\subsection*{1. Incentivising Participation}
\begin{itemize}
    \item \textbf{Staking Rewards}: Inflation can be used to reward participants who stake their tokens, helping to secure the network.
    \item \textbf{Mining Rewards}: In proof-of-work systems, inflation rewards miners for validating transactions and maintaining the blockchain.
\end{itemize}

\subsection*{2. Encouraging Network Growth}
\begin{itemize}
    \item \textbf{User Acquisition}: Newly minted tokens can be distributed to new users or developers to encourage adoption and growth of the ecosystem.
    \item \textbf{Ecosystem Development}: Inflation can fund development grants, partnerships, and other initiatives that contribute to the ecosystem's expansion.
\end{itemize}

\subsection*{3. Maintaining Network Security}
\begin{itemize}
    \item \textbf{Validator Incentives}: In proof-of-stake systems, inflation can incentivise validators to act honestly and maintain the network's security.
\end{itemize}

\subsection*{4. Economic Stability}
\begin{itemize}
    \item \textbf{Liquidity}: A controlled inflation rate can ensure there are enough tokens in circulation to facilitate transactions and maintain liquidity.
    \item \textbf{Supply Management}: Inflation can help manage the token supply, preventing issues like hoarding and ensuring tokens are available for use within the ecosystem.
\end{itemize}

\subsection*{5. Governance and Flexibility}
\begin{itemize}
    \item \textbf{Adaptive Policies}: Inflation allows for adaptive economic policies that can respond to changing network conditions and needs.
\end{itemize}

\subsection*{Example: Staking Rewards}

In a proof-of-stake blockchain, inflation might be used to reward users who lock up their tokens to help validate transactions.
This not only secures the network but also encourages long-term holding and reduces token velocity, potentially stabilising the token's value.

\subsection*{Example: Ecosystem Development}

A portion of the newly minted tokens could be allocated to a development fund, which can be used to support projects that enhance the ecosystem, such as new dApps,
infrastructure improvements, or community initiatives.

\subsection*{Balancing Inflation}

While inflation can have these benefits, it's crucial to balance it carefully to avoid negative effects such as:
\begin{itemize}
    \item \textbf{Devaluation}: Excessive inflation can devalue the token, reducing its purchasing power and attractiveness.
    \item \textbf{User Discontent}: If users feel that their holdings are being diluted too quickly, it can lead to dissatisfaction and reduced trust in the project.
\end{itemize}

Inflation, when managed properly, can be a powerful tool to incentivize participation, ensure network security, and foster ecosystem growth.
However, it requires careful planning and governance to balance the benefits against potential downsides.

\section*{Getting Your Tokens Listed and Running a Service on a Blockchain}

Getting tokens listed and running a service on a blockchain involves several steps. Here's a high-level guide to the process:

\subsection*{1. Token Creation}

\subsubsection*{Define Token Specifications}
\begin{itemize}
    \item \textbf{Blockchain Platform}: Choose a blockchain platform (e.g., Ethereum, Binance Smart Chain, Solana).
    \item \textbf{Token Standard}: Select a token standard (e.g., ERC-20 for Ethereum).
    \item \textbf{Tokenomics}: Finalize your tokenomics model (total supply, distribution, utility, etc.).
\end{itemize}

\subsubsection*{Develop Smart Contracts}
\begin{itemize}
    \item \textbf{Write Smart Contracts}: Develop the smart contracts for your token.
    \item \textbf{Audit}: Get your smart contracts audited by a reputable security firm to ensure they are secure and free of vulnerabilities.
\end{itemize}

\subsection*{2. Deploying the Token}

\subsubsection*{Deploy Smart Contracts}
\begin{itemize}
    \item \textbf{Testnet Deployment}: Deploy your smart contracts on a testnet to ensure everything works as expected.
    \item \textbf{Mainnet Deployment}: Once tested, deploy your smart contracts on the mainnet.
\end{itemize}

\subsection*{3. Running the Service}

\subsubsection*{Develop the Service}
\begin{itemize}
    \item \textbf{Backend Development}: Develop the backend infrastructure for the LokuBox service.
    \item \textbf{Frontend Development}: Create the user interface for the service.
    \item \textbf{Integration}: Integrate the LOKU token with the service, ensuring that users can interact with the blockchain as needed.
\end{itemize}

\subsubsection*{Launch the Service}
\begin{itemize}
    \item \textbf{Beta Testing}: For the full service we must conduct beta testing with a small group of users to gather feedback and make inevitable improvements.
    \item \textbf{Main Launch}: Launch the service to the public on a mainnet.
\end{itemize}

\subsection*{4. Getting Tokens Listed}

\subsubsection*{Prepare Documentation}
\begin{itemize}
    \item \textbf{Whitepaper}: There must eist a detailed whitepaper explaining the LokuBox proposition, tokenomics, and use cases.
    \item \textbf{Website}: Create a professional LokuBox focused version of the website with all relevant information about LokuBox.
\end{itemize}

\subsubsection*{Apply for Listings}
\begin{itemize}
    \item \textbf{Centralised Exchanges (CEXs)}: Apply to get your token listed on centralised exchanges like Binance, Coinbase, etc.
    This will involve filling out application forms and meeting specific requirements.
    \item \textbf{Decentralised Exchanges (DEXs)}: List your token on decentralised exchanges like Uniswap, SushiSwap, PancakeSwap, etc.
    This usually involves providing liquidity to a trading pair.
\end{itemize}

\subsubsection*{Marketing and Community Building}
\begin{itemize}
    \item \textbf{Community Engagement}: Build and engage with your community on platforms like Telegram, Discord, and Twitter.
    \item \textbf{Marketing Campaigns}: Run marketing campaigns to increase awareness and attract users and investors.
\end{itemize}

\subsection*{5. Ongoing Management}

\subsubsection*{Governance}
\begin{itemize}
    \item \textbf{DAO}: Perhaps set up a DAO (Decentralised Autonomous Organization) for community governance?
    \item \textbf{Voting}: maybe implement a voting mechanism for token holders to participate in decision-making?
\end{itemize}

\subsubsection*{Maintenance and Updates}
\begin{itemize}
    \item \textbf{Regular Updates}: Continuously improve the service and smart contracts based on user feedback and technological advancements.
    \item \textbf{Security Audits}: Conduct regular security audits to ensure the ongoing security of the smart contracts and infrastructure.
\end{itemize}

\subsection*{Example: Token Deployment on Ethereum}

Here's a simplified example of deploying an ERC-20 token on Ethereum:

\begin{verbatim}
pragma solidity ^0.8.0;

import "@openzeppelin/contracts/token/ERC20/ERC20.sol";

contract MyToken is ERC20 {
    constructor(uint256 initialSupply) ERC20("MyToken", "MTK") {
        _mint(msg.sender, initialSupply);
    }
}
\end{verbatim}

\subsection*{Deployment Script}

\begin{verbatim}
const { ethers } = require("hardhat");

async function main() {
    const [deployer] = await ethers.getSigners();
    console.log("Deploying contracts with the account:", deployer.address);

    const MyToken = await ethers.getContractFactory("MyToken");
    const token = await MyToken.deploy(1000000); // 1,000,000 tokens

    console.log("Token deployed to:", token.address);
}

main()
    .then(() => process.exit(0))
    .catch((error) => {
        console.error(error);
        process.exit(1);
    });
\end{verbatim}

\subsection*{Conclusion}

By following these steps, you can create, deploy, and list your token, as well as run a blockchain-based service.
Each step requires careful planning and execution to ensure the success and security of your project.

\section*{Protocols for Liquidity}

Liquidity protocols play a crucial role in ensuring that a token can be easily bought, sold, and traded, which is essential
for the health and success of a token ecosystem. Here's how liquidity protocols fit into the tokenomics model:

\subsection*{Importance of Liquidity Protocols}

\subsubsection*{1. Market Accessibility}
\begin{itemize}
    \item \textbf{Ease of Trading}: High liquidity ensures that users can easily trade the token without significant price slippage.
    \item \textbf{Exchange Listings}: Adequate liquidity is often a prerequisite for getting listed on major exchanges.
\end{itemize}

\subsubsection*{2. Price Stability}
\begin{itemize}
    \item \textbf{Reduced Volatility}: Liquidity helps stabilize the token price by allowing large trades to occur without drastically affecting the market price.
\end{itemize}

\subsubsection*{3. User Confidence}
\begin{itemize}
    \item \textbf{Trust}: Users are more likely to trust and invest in a token that has good liquidity, as it indicates a healthy and active market.
\end{itemize}

\subsection*{Common Liquidity Protocols}

\subsubsection*{1. Automated Market Makers (AMMs)}
\begin{itemize}
    \item \textbf{Uniswap}: A popular decentralised exchange (DEX) on Ethereum that uses liquidity pools to facilitate trading.
    \item \textbf{SushiSwap}: Another DEX that operates similarly to Uniswap but offers additional features like yield farming.
    \item \textbf{PancakeSwap}: A DEX on Binance Smart Chain (BSC) that also uses AMMs.
\end{itemize}

\subsubsection*{2. Liquidity Mining}
\begin{itemize}
    \item \textbf{Incentives}: Offering rewards (in the form of your token or another token) to users who provide liquidity to your pools.
    \item \textbf{Yield Farming}: Users can earn additional rewards by staking their liquidity provider (LP) tokens in yield farming contracts.
\end{itemize}

\subsubsection*{3. Centralised Exchange Liquidity Programs}
\begin{itemize}
    \item \textbf{Market Making}: Partnering with market makers to ensure there is always buy and sell orders for your token on centralised exchanges.
    \item \textbf{Liquidity Provision}: Some centralised exchanges offer programs where they provide liquidity in exchange for a fee or a share of the token supply.
\end{itemize}

\subsection*{An Example: Adding Liquidity on Uniswap}

Here's a simplified example of we might add liquidity to a Uniswap pool:

\subsubsection*{Step 1: Create a Liquidity Pool}
\begin{itemize}
    \item \textbf{Deploy Your Token}: Ensure your ERC-20 token is deployed on the Ethereum mainnet.
    \item \textbf{Add Liquidity}: Use the Uniswap interface or a smart contract to add liquidity.
\end{itemize}

\subsubsection*{Step 2: Add Liquidity via Smart Contract}

\begin{verbatim}
const { ethers } = require("hardhat");

async function main() {
    const [deployer] = await ethers.getSigners();
    const tokenAddress = "0xYourTokenAddress"; // Replace with your token address
    const uniswapRouterAddress = "0xUniswapV2Router02Address"; // Uniswap V2 Router address

    const Token = await ethers.getContractFactory("MyToken");
    const token = Token.attach(tokenAddress);

    const UniswapV2Router02 = await ethers.getContractAt("IUniswapV2Router02", uniswapRouterAddress);

    const amountTokenDesired = ethers.utils.parseUnits("1000", 18); // 1000 tokens
    const amountETHDesired = ethers.utils.parseEther("10"); // 10 ETH

    await token.approve(uniswapRouterAddress, amountTokenDesired);

    await UniswapV2Router02.addLiquidityETH(
        tokenAddress,
        amountTokenDesired,
        0,
        0,
        deployer.address,
        Math.floor(Date.now() / 1000) + 60 * 10, // 10 minutes deadline
        { value: amountETHDesired }
    );

    console.log("Liquidity added to Uniswap pool");
}

main()
    .then(() => process.exit(0))
    .catch((error) => {
        console.error(error);
        process.exit(1);
    });
\end{verbatim}

\subsection*{Conclusion}

Incorporating liquidity protocols into the tokenomics model is essential for ensuring market accessibility, price stability, and user confidence.
By leveraging AMMs, liquidity mining, and centralised exchange programs, you can create a robust and liquid market for your token, which is crucial for the
long-term success of the LokuBox blockchain service.

\section*{Organizing Tokens and Charging Points with zk-Rollups}

Given our service's requirements and it's potential use of zk-rollups, here's a detailed approach to organizing the LOKU tokens and charging points, along with
comments on zk-rollups and other relevant considerations:

\subsection*{Token Organization and Charging Points}

\begin{itemize}
    \item \textbf{Token Utility}:
    \begin{itemize}
        \item \textbf{Service Access}: Users need to hold and spend the LOKU token to access the file tracking service.
        \item \textbf{Transaction Fees}: Tokens are used to pay for recording changes on the blockchain.
        \item \textbf{Audit Access}: Tokens are required to access and use the audit data for forensic purposes.
    \end{itemize}
    \item \textbf{Charging Points}:
    \begin{itemize}
        \item \textbf{File Tracking}: Charge tokens for each file tracked (new, changed, viewed).
        \item \textbf{Batch Processing}: Offer discounted rates for batching multiple file changes into a single transaction.
        \item \textbf{Audit Data Access}: Charge tokens for accessing historical audit data.
    \end{itemize}
\end{itemize}

\subsection*{Strawman Token Flow}

\begin{enumerate}
    \item \textbf{User Purchases Tokens}: Users buy LOKU tokens from an exchange or directly from our platform.
    \item \textbf{Service Usage}:
    \begin{itemize}
        \item \textbf{File Tracking}: Users spend tokens to record file changes.
        \item \textbf{Batch Processing}: Users can opt to batch file changes, spending fewer tokens per file.
    \end{itemize}
    \item \textbf{Audit Access}: Users spend tokens to access and analyze audit data.
\end{enumerate}

\subsection*{zk-Rollups}

zk-Rollups can significantly enhance the efficiency and scalability of our service by batching multiple transactions into a single transaction on the main blockchain.
Here's how they can be integrated:

\begin{itemize}
    \item \textbf{Batching Transactions}:
    \begin{itemize}
        \item \textbf{Efficiency}: By batching multiple file changes into a single zk-rollup transaction,
        the number of transactions that need to be recorded on the main blockchain is reduced, lowering costs.
        \item \textbf{Scalability}: zk-Rollups handle a large number of transactions off-chain, improving the scalability of the service.
    \end{itemize}
    \item \textbf{Cost Reduction}:
    \begin{itemize}
        \item \textbf{Lower Fees}: zk-Rollups reduce the gas fees associated with recording each individual file change on the blockchain.
        \item \textbf{Economies of Scale}: The more transactions you batch, the greater the cost savings.
    \end{itemize}
    \item \textbf{Security}:
    \begin{itemize}
        \item \textbf{Zero-Knowledge Proofs}: zk-Rollups use zero-knowledge proofs to ensure the integrity and validity of the batched transactions,
        maintaining security while improving efficiency.
    \end{itemize}
\end{itemize}

\subsection*{Implementation Steps}

\begin{enumerate}
    \item \textbf{Token Smart Contract}:
    \begin{itemize}
        \item Develop and deploy LOKU ERC-20 token smart contract.
        \item Implement functions for spending tokens on file tracking and audit access.
    \end{itemize}
    \item \textbf{zk-Rollup Integration}:
    \begin{itemize}
        \item Choose a zk-rollup solution (e.g., zkSync, Loopring).
        \item Integrate the zk-rollup solution with our service to batch file changes.
    \end{itemize}
    \item \textbf{Service Backend}:
    \begin{itemize}
        \item We have the backend infrastructure to handle file tracking, batching, and audit data access.
        \item Implement logic to interact with the zk-rollup solution for batching transactions, this would replace substrate OCW.
    \end{itemize}
\end{enumerate}

\subsection*{Example: Token Smart Contract}

\begin{verbatim}
pragma solidity ^0.8.0;

import "@openzeppelin/contracts/token/ERC20/ERC20.sol";

contract ServiceToken is ERC20 {
    constructor(uint256 initialSupply) ERC20("ServiceToken", "STK") {
        _mint(msg.sender, initialSupply);
    }

    function spendTokens(address from, uint256 amount) external {
        _transfer(from, address(this), amount);
    }
}
\end{verbatim}

\subsection*{Example: zk-Rollup Integration}

\subsubsection*{Step 1: Choose a zk-Rollup Solution}
\begin{itemize}
    \item \textbf{zkSync}: A popular zk-rollup solution for Ethereum.
\end{itemize}

\subsubsection*{Step 2: Integrate zkSync}

\begin{verbatim}
const { ethers } = require("hardhat");
const zksync = require("zksync");

async function main() {
    const syncProvider = await zksync.getDefaultProvider("mainnet");
    const ethProvider = ethers.getDefaultProvider("mainnet");

    const wallet = new zksync.Wallet("YOUR_PRIVATE_KEY", syncProvider, ethProvider);

    const batch = [
        { to: "0xRecipientAddress1", token: "STK", amount: ethers.utils.parseUnits("10", 18) },
        { to: "0xRecipientAddress2", token: "STK", amount: ethers.utils.parseUnits("20", 18) },
        // Add more transactions as needed
    ];

    const batchTransaction = await wallet.batchTransfer(batch);
    console.log("Batch transaction hash:", batchTransaction.hash);
}

main()
    .then(() => process.exit(0))
    .catch((error) => {
        console.error(error);
        process.exit(1);
    });
\end{verbatim}

\subsection*{Conclusion}

By organizing our tokens and charging points effectively, and leveraging zk-rollups for batching transactions, we can create a scalable, cost-efficient,
and secure file tracking service. This approach not only reduces costs but also enhances the user experience by providing a seamless and efficient service.

\section*{Allowing Users to Buy Tokens Directly from Your Platform}

To allow users to buy tokens directly from our platform, we must implement a token sale mechanism, often referred to as a "crowdsale" or "initial coin offering (ICO)."
Here's a step-by-step guide to setting this up:

\subsection*{1. Smart Contract for Token Sale}

We need a smart contract that handles the token sale. This contract will accept payments (usually in ETH) and distribute your tokens to buyers.

\subsubsection*{Example: Token Sale Smart Contract}

\begin{verbatim}
pragma solidity ^0.8.0;

import "@openzeppelin/contracts/token/ERC20/IERC20.sol";
import "@openzeppelin/contracts/access/Ownable.sol";

contract TokenSale is Ownable {
    IERC20 public token;
    uint256 public rate; // Number of tokens per ETH
    uint256 public startTime;
    uint256 public endTime;

    event TokensPurchased(address indexed purchaser, uint256 amount);

    constructor(
        IERC20 _token,
        uint256 _rate,
        uint256 _startTime,
        uint256 _endTime
    ) {
        require(_rate > 0, "Rate must be greater than 0");
        require(_startTime < _endTime, "Start time must be before end time");

        token = _token;
        rate = _rate;
        startTime = _startTime;
        endTime = _endTime;
    }

    modifier onlyWhileOpen() {
        require(block.timestamp >= startTime && block.timestamp <= endTime, "Sale is not open");
        _;
    }

    function buyTokens() public payable onlyWhileOpen {
        uint256 tokenAmount = msg.value * rate;
        require(token.balanceOf(address(this)) >= tokenAmount, "Not enough tokens in contract");

        token.transfer(msg.sender, tokenAmount);
        emit TokensPurchased(msg.sender, tokenAmount);
    }

    function withdrawFunds() public onlyOwner {
        payable(owner()).transfer(address(this).balance);
    }

    function withdrawRemainingTokens() public onlyOwner {
        uint256 remainingTokens = token.balanceOf(address(this));
        token.transfer(owner(), remainingTokens);
    }
}
\end{verbatim}

\subsection*{2. Deploy the Smart Contract}

Deploy both our token contract and the token sale contract to the Ethereum mainnet (or another blockchain of choice).

\subsection*{3. Frontend Integration}

Create a user-friendly interface on the LokuBox platform where users can interact with the token sale contract.

\subsubsection*{Example: Frontend Integration with Web3.js}

\begin{enumerate}
    \item \textbf{Install Web3.js}:
    \begin{verbatim}
    npm install web3
    \end{verbatim}
    \item \textbf{Create a Simple Purchase Form}:
    \begin{verbatim}
    <!DOCTYPE html>
    <html lang="en">
    <head>
        <meta charset="UTF-8">
        <meta name="viewport" content="width=device-width, initial-scale=1.0">
        <title>Token Sale</title>
    </head>
    <body>
        <h1>Buy Tokens</h1>
        <form id="buyTokensForm">
            <label for="ethAmount">Amount of ETH:</label>
            <input type="text" id="ethAmount" name="ethAmount">
            <button type="submit">Buy Tokens</button>
        </form>

        <script src="https://cdn.jsdelivr.net/npm/web3/dist/web3.min.js"></script>
        <script src="app.js"></script>
    </body>
    </html>
    \end{verbatim}
    \item \textbf{Handle Token Purchase in JavaScript}:
    \begin{verbatim}
    const web3 = new Web3(Web3.givenProvider || "http://localhost:8545");

    const tokenSaleAddress = "0xYourTokenSaleContractAddress"; // Replace with your TokenSale contract address
    const tokenSaleABI = [ /* ABI of your TokenSale contract */ ];

    const tokenSaleContract = new web3.eth.Contract(tokenSaleABI, tokenSaleAddress);

    document.getElementById("buyTokensForm").addEventListener("submit", async (event) => {
        event.preventDefault();

        const ethAmount = document.getElementById("ethAmount").value;
        const accounts = await web3.eth.requestAccounts();
        const account = accounts[0];

        tokenSaleContract.methods.buyTokens().send({
            from: account,
            value: web3.utils.toWei(ethAmount, "ether")
        })
        .on("transactionHash", (hash) => {
            console.log("Transaction sent: ", hash);
        })
        .on("receipt", (receipt) => {
            console.log("Transaction confirmed: ", receipt);
        })
        .on("error", (error) => {
            console.error("Transaction failed: ", error);
        });
    });
    \end{verbatim}
\end{enumerate}

\subsection*{4. Security and Compliance}

\begin{itemize}
    \item \textbf{KYC/AML}: Implement Know Your Customer (KYC) and Anti-Money Laundering (AML) procedures if required by your jurisdiction.
    \item \textbf{Smart Contract Audits}: Ensure our smart contracts are audited by a reputable security firm to prevent vulnerabilities.
\end{itemize}

\subsection*{5. Marketing and Community Engagement}

\begin{itemize}
    \item \textbf{Announcements}: Announce the LOKU token sale on social media, forums, and through press releases.
    \item \textbf{Community Building}: Engage with our community on platforms like Telegram, Discord, and Twitter to build trust and excitement around LokuBox.
\end{itemize}

\subsection*{Conclusion}

By implementing a token sale smart contract and integrating it with a user-friendly frontend, we allow users to buy tokens directly from our platform.
We must ensure we follow best practices for security and compliance to build trust and protect our users.
\bibliography{tokenomics}
\bibliographystyle{plainnat}
\end{document}